% =====================================================================
\chapter*{Introduction}
\addcontentsline{toc}{chapter}{Introduction}
\markboth{Introduction}{Introduction}
\label{chap-introduction}
% =====================================================================

\section*{Contexte}
\addcontentsline{toc}{section}{Contexte}

Les grands modèles de langue ont une mémoire figée à la date de leur
entraînement et une propension bien documentée à produire des affirmations
plausibles mais fausses lorsqu'on les interroge sur des faits qu'ils n'ont pas
mémorisés. La génération augmentée par récupération~\citep{rag2020} répond à ces
deux limites par la même mécanique~: avant de répondre, le système va chercher
dans une base documentaire les passages pertinents pour la question, puis les
fournit au modèle génératif comme contexte. Le modèle n'a plus à savoir,
il a à lire. Cette architecture est devenue le mode de déploiement
dominant des modèles de langue en environnement professionnel, et une
littérature abondante en a exploré les variantes~\citep{ragsurvey2023}.

Le maillon faible de cette chaîne est rarement le générateur. C'est le
prétraitement documentaire. Un pipeline RAG classique commence par convertir
chaque document en texte brut~: analyse syntaxique, extraction du flux
textuel, reconnaissance optique de caractères sur un document
numérisé. Cette conversion est une projection destructive. Un tableau devient
une suite de mots dont on ne sait plus quelle cellule appartenait à quelle
ligne. Une légende se retrouve orpheline de sa figure. Un encadré latéral
s'intercale au milieu d'un paragraphe qu'il ne concernait pas. Cette perte d'information n'est pas marginale sur des
documents que l'on qualifie de visuellement riches comme des pages web,
des rapports financiers, des notices techniques ou des articles encyclopédiques. La perte d'information porte précisément sur ce que la mise en page encodait.

Deux réponses distinctes coexistent dans la littérature récente, et ce
travail ci se situe exactement à leur intersection. La première rend
l'extraction inutile~: puisque les modèles vision-langage savent lire une
image de document, autant leur donner l'image, une page rendue, découpée en
imagettes indexées par un encodeur multimodal, la récupération s'opérant
directement dans l'espace visuel des pixels~\citep{dse2024, colpali2024}.
PixelRAG~\citep{pixelrag2026} porte cette idée à l'échelle du web et rapporte
qu'elle égale ou dépasse la récupération textuelle sur les documents longs et
fortement mis en page~; mais elle découpe la page rendue sur une grille de
hauteur fixe, indifférente à ce qu'elle traverse. La seconde structure les
preuves plutôt que de les traiter comme un sac plat d'unités indépendantes~:
MAGE-RAG~\citep{magerag2026} organise un document déjà segmenté en un graphe
multigranulaire qu'un contrôleur budgété parcourt à la requête, plutôt que de présenter au modèle lecteur la totalité du document ou une sélection figée.

Ces deux travaux sont complémentaires et laissent la même
pièce sur la table. PixelRAG établit que des tuiles de pixels lues par un
modèle vision-langage constituent une unité de récupération viable, mais les
découpe à l'aveugle. MAGE-RAG établit qu'un contrôleur de graphe budgété
surpasse une sélection top-$k$ statique, mais suppose les éléments du document
déjà constitués, en amont de son propre pipeline. Ce travail assemble les deux
par l'artefact que PixelRAG jette et que MAGE-RAG n'a jamais eu à construire~:
le DOM de la page rendue.

\section*{Le problème}
\addcontentsline{toc}{section}{Le problème}

Ce changement de représentation laisse intacte une décision de conception
sur la manière dont on doit découper la page. Un modèle vision-langage
ne peut pas ingérer une capture d'écran pleine page d'un article
encyclopédique de \num{46000} pixels de hauteur, et de toute façon une unité de
récupération doit être suffisamment fine pour qu'un score de pertinence
signifie quelque chose. Il faut donc découper la page en plusieurs captures. PixelRAG découpe sur une
grille de hauteur fixe, sans aucune notion de ce que la découpe
traverse. Chaque bande a la même hauteur, qu'elle contienne un paragraphe
entier, la moitié d'un tableau, ou le bas d'une image et le haut de la section
suivante.

La conséquence est mécanique. Une réponse courte noyée dans une bande de \num{1024} pixels qui agrège trois
sections hétérogènes devient très difficile à retrouver. De ce fait, le signal utile
représente quelques mots sur plusieurs centaines. Symétriquement, un tableau
coupé au milieu produit deux unités dont aucune ne porte l'information
complète. Le système de récupération se voit alors confier une tâche
supplémentaire qui n'était pas la sienne~: reconstruire, à partir des pixels,
les frontières sémantiques du document.

L'observation qui fonde ce travail est que cette reconstruction est inutile pour une page web. Le moteur de mise en page du navigateur a déjà
calculé la boîte englobante exacte de chaque élément avant qu'un seul pixel ne
soit peint. La segmentation que le RAG visuel dépense de l'effort à
reconstruire est disponible, exacte, et gratuite, dans le DOM de la page
rendue. Il ne s'agit ni d'un modèle à entraîner, ni d'une heuristique de
segmentation visuelle à calibrer, ni d'un détecteur de frontières appris~: il
s'agit d'une information que le navigateur expose par une simple requête.
La \cref{fig-intro_teaser} montre directement ce que cela change sur une page
réelle.

\begin{figure}[t]
\centering
\begin{subfigure}[b]{0.545\textwidth}
  \centering
  \includegraphics[trim=0 0 0 22pt, clip, width=\linewidth]{tiling_overlay.pdf}
  \caption{Tuilage guidé par le DOM}
  \label{fig-intro_teaser_dom}
\end{subfigure}
\hfill
\begin{subfigure}[b]{0.425\textwidth}
  \centering
  \includegraphics[width=\linewidth]{grid_overlay_illustrative.pdf}
  \caption{Grille de hauteur fixe}
  \label{fig-intro_teaser_grid}
\end{subfigure}
\caption[La même page rendue, tuilée de deux manières]{La même page Wikipédia
rendue (\code{en.wikipedia.org/wiki/George\_P.\_Gunn}), découpée de deux
manières. \textbf{(a)}~Tuilage guidé par le DOM (\cref{sec-tuilage})~:
\textbf{(b)}~Grille aveugle de hauteur fixe selon PixelRAG}
\label{fig-intro_teaser}
\end{figure}

\section*{Contribution personnelle}
\addcontentsline{toc}{section}{Contribution personnelle}

L'intégralité du système décrit dans ce rapport, avec la conception du tuilage,
la construction du graphe, l'entraînement contrastif du lecteur, le contrôleur
budgété, ainsi que le code, les jeux de données synthétiques et les
évaluations qui les accompagnent ont été conçues et implémentées par l'auteur
dans le cadre de ce projet de master. Les deux seuls emprunts à des travaux
tiers, explicitement identifiés comme tels tout au long du rapport, sont
l'architecture de graphe et de contrôleur de MAGE-RAG~\citep{magerag2026},
dont ce travail généralise le contrôleur à des frontières inter-documents
qu'il ne traite pas, et le protocole de grille de PixelRAG~\citep{pixelrag2026},
reproduit à l'identique dans un module indépendant pour servir de condition de
référence à l'\cref{sec-res_ablation_a}. Tout le reste est une contribution
originale de ce projet.

Une version antérieure et condensée des \cref{chap-methodologie,chap-resultats}
a été soumise en double aveugle, sous le titre \emph{\og DOM-Graph RAG:
DOM-Guided Adaptive Tiling and Budgeted Online Evidence Control for
Multimodal Web-Page RAG\fg{}}~\citep{domgraphrag2026}, à l'atelier
\emph{DocInsights 2026}, associé à EMNLP~2026. Le présent rapport en constitue
la version longue~: il ajoute le contexte scientifique complet
(\cref{chap-etat_de_lart}), le détail d'implémentation et d'ingénierie
nécessaire à la reproduction, et l'analyse critique
des quatre évaluations du \cref{chap-resultats}, dont la version courte ne
pouvait rapporter que les résultats les plus saillants.

\section*{Objectifs et contributions}
\addcontentsline{toc}{section}{Objectifs et contributions}

L'objectif de ce projet de master est de transformer l'observation ci-dessus en
un système complet et de mesurer ce qu'elle apporte réellement. Cela suppose de construire quatre composants et de
les évaluer séparément puis ensemble. Les contributions de ce travail sont les
suivantes.

\begin{enumerate}
\item \textbf{Un tuilage adaptatif guidé par le DOM} (\cref{sec-tuilage}) qui
remplace la grille de pixels aveugle de PixelRAG~: chaque élément est recadré
sur sa propre boîte englobante, les éléments trop petits sont fusionnés avec
leurs voisins, ceux trop grands sont subdivisés, et les titres sont
systématiquement appariés avec leur contenu. Le point technique le plus
délicat est la mesure de la boîte englobante des balises de
texte coulant, dont la boîte CSS s'étend sur toute la largeur du conteneur
même lorsque le texte s'enroule autour d'un élément flottant. C'est cette
étape, formalisée en détail au \cref{chap-methodologie}, qui porte
la recherche centrale de ce travail~: aligner la grille de \anglais{patches}
que consomme le modèle vision-langage sur les frontières sémantiques du DOM
plutôt que sur une géométrie arbitraire, de sorte que chaque unité visuelle
soumise à l'encodeur porte un contenu homogène et complet.

\item \textbf{Un graphe de preuves multigranulaire} (\cref{sec-graphe})
organisant les tuiles en deux niveaux (page et élément) et instanciant les cinq
relations du graphe de MAGE-RAG~\citep{magerag2026}, augmentées de deux
relations propres à notre contexte~: \code{links\_to}, qui suit les liens
hypertextes pour relier plusieurs pages, et \code{continues}, qui chaîne les
sections d'une page trop longue pour être traitée d'un bloc.

\item \textbf{Un lecteur vision-langage entraîné par contraste}
(\cref{sec-contrastif}) sur des paires (question, tuile) générées
synthétiquement avec négatifs difficiles minés, adapté par LoRA sur le décodeur
et la tour de vision, via une tête d'encodage de type
PromptEOL~\citep{prompteol2024}.

\item \textbf{Un contrôleur de preuves en ligne budgété} (\cref{sec-controleur})
qui fait évoluer l'état de chaque nœud du graphe selon une politique
best-first sous budget, généralisant l'algorithme de contrôle de MAGE-RAG à des
frontières qui franchissent les limites de document.
\end{enumerate}

À ces contributions logicielles s'ajoute une contribution expérimentale~:
\textbf{quatre évaluations} (\cref{chap-resultats}) qui isolent chaque
composant puis mesurent le système assemblé, avec une attention particulière
portée à la comparabilité des métriques. Nous montrons notamment que le
Recall@$k$ brut n'est pas comparable entre conditions de tuilage, puisqu'un
vivier de candidats plus petit le gonfle mécaniquement, et nous utilisons à la
place un rang percentile invariant à la taille du vivier.

\section*{Résultats principaux}
\addcontentsline{toc}{section}{Résultats principaux}

Deux résultats sont nets. Le tuilage guidé par le DOM double presque la
proportion de questions dont la réponse survit dans une unité récupérable
(\SI{51.8}{\percent} contre \SI{26.4}{\percent} pour la grille de PixelRAG),
tout en consommant \num{2.7} fois moins de pixels par page~; ces deux effets
vont dans le même sens, ce qui est contre-intuitif puisque notre tuilage
produit cinq fois plus de tuiles. L'entraînement contrastif du lecteur fait
presque doubler le Recall@1 sur un partitionnement par article strictement
disjoint ($\num{0.433} \rightarrow \num{0.785}$).

Les composants en aval donnent des résultats plus nuancés, que nous analysons
en détail plutôt que de les minimiser. L'avantage de classement du tuilage DOM
sur la grille est réel mais modeste, et il s'inverse face à un découpage
textuel plat sur le sous-ensemble strictement apparié~: nous montrons que ce
résultat est cohérent avec le scoreur lexical employé, qui n'exploite aucune
information spatiale, et non avec un défaut du tuilage lui-même. Le contrôleur
budgété n'établit pas encore de supériorité sur une sélection top-$k$
statique~; nous en diagnostiquons une cause architecturale identifiable, la
corrigeons, mesurons le gain, et documentons ce qui reste à faire. Enfin,
l'évaluation bout-en-bout montre que le système assemblé échoue par
abstention plutôt que par fabulation, ce qui désigne le sous-rappel de
preuves comme le verrou restant et pointe vers un correctif précis et
déjà instrumenté dans le code.

\section*{Organisation du rapport}
\addcontentsline{toc}{section}{Organisation du rapport}

Le rapport suit une progression en trois chapitres.

Le \textbf{\cref{chap-etat_de_lart}} situe brièvement DOM-Graph RAG dans son
contexte scientifique avec la formulation du problème de récupération et de ses
métriques, le prétraitement des documents visuellement riches, les modèles
vision-langage et la famille du RAG visuel, le RAG structuré par graphe, et
les techniques d'apprentissage contrastif mobilisées ensuite. Il reste volontairement
concis et se clôt sur l'identification précise du
verrou et le positionnement de ce travail face à PixelRAG et MAGE-RAG.

Le \textbf{\cref{chap-methodologie}} constitue le cœur de ce rapport. Il ouvre
sur ce que le navigateur expose réellement de la géométrie d'une page rendue puis
suit l'ordre du pipeline~: rendu et extraction géométrique, tuilage adaptatif
formalisé en détail, construction du graphe, entraînement contrastif du
lecteur et contrôleur en ligne. Chaque étape est explicitée par ses choix de
conception, sa formalisation mathématique, les cas limites rencontrés sur des
pages réelles, et les décisions d'ingénierie qui rendent le pipeline exécutable
sur du matériel modeste.

Le \textbf{\cref{chap-resultats}} présente les quatre évaluations, leur
analyse critique, les coûts de calcul observés, et une discussion structurée des
limites du travail et des directions qu'elles ouvrent.

Un chapitre de conclusion synthétise l'apport et les enseignements du projet,
et les annexes rassemblent les hyperparamètres, l'organisation du dépôt, les
invites utilisées et les protocoles de reproduction.
